\chapter{能量下降还不够：熵与自由能}

\begin{kaichang}
体育课上有人扭了脚，老师从医药箱里摸出一只扁扁的塑料袋——一次性速冷冰袋。她两手一捏，袋子里“啪”地一声轻响，几秒钟后，整只袋子变得冰凉，敷在脚踝上直冒寒气。

现在请你当一回侦探。这只袋子变冷，说明里面的过程在\textbf{吸热}——它正从你的手、从空气里把热量抢走。可是没有任何人给它加热、点火、通电，它就这么自己发生了。再看桌上：一块冰在室温里慢慢融化，同样吸热，同样自己发生，谁也拦不住。

这就怪了。第 27 章我们刚立下一条规矩：变化愿意朝“能量降低”的方向走，放热的过程才有“动力”。可冰袋和冰块明明是反着来的——它们吸着热，却照样自发进行。是规矩错了，还是规矩漏掉了什么？先别翻答案，猜猜看：如果能量下降还不够，那另一个裁判会是谁？
\end{kaichang}

\section{被打破的能量规矩}

先把能看见的事实摆出来。一次性冰袋里装的是硝酸铵（\chem{NH_4NO_3}）固体和一个水囊，捏破水囊，硝酸铵开始溶解。这个溶解过程要吸热：每溶解 \ud{1}{mol} 硝酸铵，大约从环境夺走 \ud{26}{kJ} 的热量——所以袋子变冷。用第 27 章的语言说，这是一个 $\Delta H>0$ 的过程，体系的能量账是“亏”的。按“能量越低越愿意走”的老规矩，它根本不该自发发生。可它偏偏发生了，而且一旦发生就拦不住，除非你把它塞回冰箱冻上好久。

冰袋不是个例，生活里这样的“违规者”一抓一大把：

\begin{itemize}[itemsep=2pt]
\item 冰块在 $25\,^{\circ}\mathrm{C}$ 的房间里融化，吸热，自发。
\item 打开香水瓶盖，不用摇不用吹，香味自己扩散到整个房间——气体分子从聚集的一小团“摊”成大空间里的稀薄分布，同样没有能量下降的驱动。
\item 一勺糖丢进热水，即使完全不搅拌，过上足够久，整杯水也会一样甜。
\item 拧开碳酸饮料的瓶盖，溶解的 \chem{CO_2} 争先恐后往外跑——跑出去要吸热，它们还是跑了。
\end{itemize}

当然也有“守规矩”的：热水放在桌上会慢慢凉下来，热量从集中的热水摊向整个房间，能量账和方向一致。麻烦的是这些不守规矩的。它们提示我们：焓变 $\Delta H$ 是裁判之一，但不是唯一的裁判。自然界手里还有第二本账。

\begin{shiyan}
\textbf{给“吸热却自发”称一称重。}材料：一只市售一次性速冷冰袋（药店或网店有售，成分是硝酸铵和水）、一支室温温度计、一条干毛巾、一个计时器。步骤：① 先记下室温，用手背摸摸袋子，感受它的初始温度；② 隔着毛巾捏破内袋，轻摇十秒，启动计时；③ 每隔 30 秒摸一次外壁（温度计贴在外壁上读数更好），记录温度何时最低、之后怎样回升；④ 同时观察：袋子外面有没有结冰？步骤④是理解本章的关键一问。安全提示：全过程\textbf{不要拆开外袋}，不要尝或挤里面的东西到皮肤上，用完包好丢弃；皮肤敏感者全程隔毛巾操作。观察点：它自发进行、自发吸热、却始终没有结冰——记住这三个“却”，本章结尾你要能逐条解释。如果没有冰袋，替代观察：在杯子里放冰块，记录室温下融化全过程的温度变化，同样问自己“它吸热，为什么不用人帮忙”。
\end{shiyan}

\section{粒子为什么爱“摊开”}

要找到第二本账，得把镜头推回粒子层。先问一个看起来很天真的问题：香味为什么会扩散？分子又没有脚，没有意愿，凭什么“非摊开了不可”？

答案是：它们什么也不凭，纯粹是概率。想象一个盒子，中间用隔板分成左右两半，左边放 4 个气体分子，右边真空。抽掉隔板后，每个分子在左边或右边的可能性各占一半，就像抛硬币。4 个分子全部恰好待在左边的排法只有 1 种；而“两左两右”的排法有 6 种——你不妨自己数一数：甲乙、甲丙、甲丁、乙丙、乙丁、丙丁。分子越多，悬殊越吓人：

\begin{qiaoliang}
把 $N$ 个分子放进左右两室的盒子，每个分子独立地随机选一边。全部 $N$ 个分子恰好挤在同一边的排法只有 1 种，而总排法有 $2^N$ 种，所以“全在一边”的概率是 $1/2^N$。

$N=4$：概率 $1/16$，偶尔还能撞见。$N=100$：概率约 $1/10^{30}$——就算你每秒观察一次，从宇宙大爆炸看到今天（约 $4\times10^{17}$ 秒），也连一次都等不到。而你房间里约有 $10^{27}$ 个空气分子，它们全部自发挤回半个房间的概率是 $1/2^{10^{27}}$，这个数小到连写出来都做不到。

反过来说，“分子摊满整个盒子”并不是分子追求的目标，而是因为它对应着\textbf{压倒性多数的排法}。均匀分布不是一个特殊状态，而是亿亿亿……种几乎无法区分的状态的总称。
\end{qiaoliang}

\begin{center}
\begin{tikzpicture}[scale=0.9]
\draw (0,0) rectangle (4,2);
\draw[dashed] (2,0) -- (2,2);
\fill (0.6,0.5) circle (2pt); \fill (1.2,1.4) circle (2pt);
\fill (0.9,1.0) circle (2pt); \fill (1.5,0.4) circle (2pt);
\fill (0.4,1.6) circle (2pt); \fill (1.4,1.7) circle (2pt);
\node at (2,-0.4) {抽掉隔板前：排法很少};
\draw (5.5,0) rectangle (9.5,2);
\fill (6.1,0.5) circle (2pt); \fill (8.7,1.4) circle (2pt);
\fill (6.9,1.0) circle (2pt); \fill (8.4,0.4) circle (2pt);
\fill (6.4,1.6) circle (2pt); \fill (8.9,1.7) circle (2pt);
\node at (7.5,-0.4) {之后：摊开的排法极多};
\end{tikzpicture}
\end{center}

（这是人为的表示法：圆点大小不代表真实分子，分子也不会停在原地。）

能量也在玩同一个游戏。热水里的分子带着“集中”的热运动能量，桌上空气分子带着较少的能量；两者碰撞来碰撞去，能量有许多许多种“摊匀”的分法，却只有极少几种“越来越集中”的分法。于是热量总是从热流向冷——不是规律禁止反向，而是反向的概率小到等于零。

把两节内容合起来看，“摊开”其实有三副面孔：粒子在\textbf{空间}里摊开（香味扩散）、能量在粒子\textbf{之间}摊开（热水变凉）、不同粒子互相\textbf{混合}着摊开（糖溶解、盐水和纯水连通后变匀）。三副面孔背后是同一笔账：可选的排法变多了。这个视角有个朴素的名字：热力学第二定律——孤立体系的熵总是趋向增大，因为“多”永远压倒“少”。你也可以反着想——你见过胡椒粉和盐摇匀之后，再摇一摇自动分开吗？见过放凉的咖啡自己重新热起来、顺便把桌子冻出霜吗？这些“倒放”场景不违反任何力学定律，却永远不会发生，因为它们要求亿万个粒子同时挤进极少数的排法里。

现在可以给第二本账起名字了。一种宏观状态背后对应的微观排法越多，我们就说它的\textbf{熵}越大。熵就是给“摊法”计数的账：粒子摊得越开、能量摊得越匀，排法越多，熵越高。

\section{熵：数一数有多少种摊法}

把上面的直觉钉成符号。熵记作 $S$，它的变化记作 $\Delta S$。因为“摊法”的数目跟粒子数、能量多少有关，熵的单位是能量除以温度：$\mathrm{J/K}$（焦每开尔文）。数排法的工作听起来抽象，但有几条朴素推论你立刻能用：

\begin{itemize}[itemsep=2pt]
\item 同样的物质，\textbf{气体的熵远大于液体，液体大于固体}。气体分子在整个容器里乱窜，排法天文数字；固体里粒子被锁在晶格上（第 24 章），只能在原地发抖，排法少得可怜。
\item 温度越高，粒子能量越分散，熵越大。
\item 所以融化（$\Delta S>0$）、蒸发（$\Delta S\gg 0$）、固体溶解、气体分子数变多的反应，都是熵增；反过来，结冰、凝结、结晶是体系的熵减。
\end{itemize}

请务必注意措辞：熵是“\textbf{微观排法的数目}”，不是“乱”。整齐的晶体熵低，不是因为格子“整洁好看”，而是因为粒子位置被定死了、可选的排法少；混乱这个词只是日常联想，不是定义。把熵当成“计数”而不是“脏乱”，后面的账才不会算错。

再练两个例子找感觉。碳酸饮料冒泡：\chem{CO_2} 从被水分子团团围住的溶解状态冲进空气，可选的空间暴涨，$\Delta S>0$——哪怕逸出过程吸点热，熵账也常常是赚的。给气体加热：分子平均动能变大，每个分子能取的能量“档位”变多，排法随之变多，熵也增大。反过来，把热汤盖上盖子焖着，并没有改变汤的熵本身，只是拦住了它向环境摊熵的通道——这也解释了为什么“熵增”总跟“拦不住的扩散、传热、混合”结伴出现。

\begin{wuqu}
“熵就是混乱程度，熵增就是东西越来越乱；房间不收拾熵就增加。”——这个比喻半对半错，错的那半会害人。反例：把一杯过饱和的醋酸钠溶液（暖宝宝里那种）静置，轻轻一晃，针状的晶体唰唰地长满整杯——粒子从自由游荡变成排列极其整齐的晶格，体系的“混乱程度”明明下降了，这个过程却照样自发！原因很简单：结晶放热，热量摊给环境，环境的熵增加得比晶体减少的还多，\textbf{体系加环境的总账}仍是熵增。生命更是如此：你的身体每秒都在把散乱的分子组装成高度有序的细胞，从不违反任何定律，因为你同时通过呼吸、散热向环境排出大量的熵。只盯着体系的“整齐不整齐”，而不算环境的账，就会得出“生命和结晶违反热力学”的荒唐结论。
\end{wuqu}

\section{环境也有发言权：自由能}

现在可以破案了。冰袋里的硝酸铵溶解：体系吸热，$\Delta H>0$，能量账吃亏；但溶解后离子从锁死的晶格里解放出来，在整袋水里自由游荡（第 25 章的水合离子），熵大幅增加，$\Delta S>0$。一本账亏，一本账赚——谁说了算？

十九世纪的物理学家发现，恒温恒压下（也就是你日常生活的条件），可以把两本账合并成一个数字。关键在于：体系放出的热量不会消失，它被环境吸收，增加了环境的熵；而且同样一份热量，在低温环境里引起的熵增更大（冷环境本来“平静”，多一点热量扰动明显；热环境本来就沸腾，同样热量只是杯水车薪）。所以环境的熵账由 $-\Delta H/T$ 决定。体系自己的熵账是 $\Delta S$。总账为正，过程就能自发。把总账整理成一个属于体系的量，就是美国物理学家吉布斯提出的\textbf{吉布斯自由能} $G$：

\[\Delta G = \Delta H - T\Delta S\]

读法只有一句：$\Delta G<0$，过程在恒温恒压下自发；$\Delta G>0$，反方向自发；$\Delta G=0$，两边扯平，达到平衡。焓账 $\Delta H$ 是“能量想往低处走”的老裁判，$T\Delta S$ 是新裁判“排法想往多处走”的分量，温度 $T$ 是两者的汇率——温度越高，熵那本账越值钱。四种组合，一目了然：

\begin{table}[h]
\centering
\begin{tabular}{cccl}
\toprule
$\Delta H$ & $\Delta S$ & $\Delta G$ & 结论与例子 \\
\midrule
负（放热） & 正（熵增） & 恒为负 & 永远自发，如燃烧 \\
正（吸热） & 正（熵增） & 高温转负 & 升温才自发，如冰袋、石灰石煅烧 \\
负（放热） & 负（熵减） & 低温转负 & 降温才自发，如水结冰 \\
正（吸热） & 负（熵减） & 恒为正 & 永远不自发 \\
\bottomrule
\end{tabular}
\end{table}

表格第二、三行藏着本章标题的全部含义：\textbf{能量下降还不够}。吸热过程只要熵增够大、温度够高，$\Delta G$ 照样为负——冰袋、融冰、碳酸饮料冒泡，全在这一行。反过来，水结冰放热，冬天能自发，夏天却不行，因为 $T$ 变大后 $-T\Delta S$ 项翻了盘。这就是为什么\textbf{温度能扳动变化的方向}：$0\,^{\circ}\mathrm{C}$ 是水和冰的扯平点，$\Delta G=0$；高于它，熵账赢，冰化水；低于它，焓账赢，水结冰。

工业上把这个原理用得炉火纯青。石灰窑里烧石灰石：\[\chem{CaCO_3} \rightarrow \chem{CaO} + \chem{CO_2}\] 这反应吸热，乍一看“不该发生”；但它放出气体，熵猛增，只要窑温烧到约 $900\,^{\circ}\mathrm{C}$ 以上，$T\Delta S$ 压过 $\Delta H$，$\Delta G$ 转负，石灰石就源源不断地分解成生石灰。地球上的建筑能有石灰和水泥，靠的是工人把温度抬过了熵账翻盘的临界点——而降低温度，方向立刻反过来。同一条式子，也解释了为什么有些反应加热反而“变懒”：水蒸气在高温下就是不肯轻易凝结。

不放心？拿真数字算一笔冰的账。融化 \ud{1}{mol} 冰（一小勺，\ud{18}{g}）要吸热约 \ud{6.0}{kJ}。这笔热量由环境出，所以环境熵减少 $6000/T$；而水分子挣脱晶格，体系熵增加约 $6000/273\approx\ud{22}{J/K}$（用熔点估算）。在 $25\,^{\circ}\mathrm{C}$（$T=298$ K）的房间里：环境账 $-6000/298\approx-\ud{20}{J/K}$，总账 $22-20=+\ud{2}{J/K}$，大于零——冰自发融化。换到 $-10\,^{\circ}\mathrm{C}$（$T=263$ K）的冷冻室：环境账变成 $-6000/263\approx-\ud{23}{J/K}$，总账 $22-23=-\ud{1}{J/K}$，小于零——冰不化，反倒是水会自发结冰。两个方向之间只隔一个温度：$T=273$ K，也就是 $0\,^{\circ}\mathrm{C}$，恰好总账为零。你看，“冰在零度以上融化”这条从小背到大的常识，原来是一道精确的算术题。

以后遇到任何过程，可以照这个清单走三步：第一步，问焓账——成键多还是断键多，$\Delta H$ 是正还是负；第二步，问熵账——粒子是更自由了（气体变多、固体溶解、温度升高）还是更受拘束，$\Delta S$ 是正还是负；第三步，看温度——把两本账按 $T$ 折合成 $\Delta G$。三步走完，方向就清楚了；至于走多快、要不要推一把，那是清单管不了的事。

\section{自发不等于快速}

先别急着把 $\Delta G<0$ 当成“马上发生”的同义词。看三个事实：

\begin{itemize}[itemsep=2pt]
\item 你手上钻戒（如果有的话）的金刚石，在常温常压下 $\Delta G<0$——它“应该”变成石墨。可它放上几亿年也纹丝不动，珠宝商从不为此失眠。
\item 铁钉生锈，$\Delta G$ 显然是负的，但一根钉子锈穿要好几年；同样是氧化，炸药爆炸只要百万分之一秒。
\item 氢气和氧气混在一个瓶子里，生成水的 $\Delta G$ 负得惊人，可它们能和平共处一整年，什么都不发生——直到你递进去一个火花。
\end{itemize}

这些例子逼我们区分两个问题：\textbf{“朝不朝这个方向走”}和\textbf{“走得有多快”}。$\Delta G$ 只回答第一个。它告诉你终点线在哪边，却不管路上的坡有多高。金刚石变石墨，方向是通的，但每个碳原子都要挣脱三根牢固的共价键（第 24 章）才能换个位置重排——这道坎高得离谱，室温下几乎没有碳原子攒得够力气翻过去。于是“自发”的过程被卡在半路上，一卡就是亿年。

反过来，同样 $\Delta G<0$，不同过程的“脚程”天差地别：冰袋几秒就变凉，暖宝宝拆开要几分钟才热透（它也是铁氧化的放热反应），一根栏杆锈穿要几年，而岩石风化出土壤要以万年计。方向的判决书一模一样，执行的日程表却毫不相干——判决书管方向，日程表另有制定者。

所以两句话都要记住：自发不等于快速；自发也不等于无需触发。氢氧混合气需要火花，不是火花“提供了反应的动力”——反应放出的能量比火花那点能量多几万倍——而是火花帮第一批分子翻过那道坎，之后放出的热自会推着后面的分子接着翻。鞭炮、煤气灶、汽车发动机，全都是这个套路：方向早就写好了，缺的只是一把开门的钥匙。方向是一本账，速度是另一本账；这第二本账的主角叫活化能，那是第 29 章的正题，这里先埋个引子。

\begin{zhengju}
“熵是排法的数目”这个今天看来顺理成章的想法，当年是一场赌上学术生命的争论。十九世纪七十年代，奥地利物理学家玻尔兹曼主张：热学的一切规律都能还原为原子的统计行为。他证明，一个孤立体系的熵只可能增加或不变（热力学第二定律），原因不是某条铁律“禁止”熵减，而是熵减对应的排法太少、概率太低。他把答案浓缩成一个公式，$S=k\ln W$：熵等于一个常数乘上微观排法数 $W$ 的对数。

麻烦在于：当时没人见过原子。以马赫为首的一派权威学者认为，原子不过是计算方便的虚构，把整个热力学押在“看不见的小球”上是危险的玄学。争论激烈到什么程度？玻尔兹曼要不断写文章、做演讲为自己辩护，对手甚至拒绝承认气体分子存在。更微妙的是数学上的诘难：既然分子运动规律本身没有时间方向（倒放一段分子碰撞的影片，每一步都不违反力学），凭什么宏观世界有方向？玻尔兹曼的回答——方向来自概率，不来自力学——超越了整个时代，很多人听不懂，也不买账。

1906 年，玻尔兹曼在抑郁中去世，没能等到终局。而翻案的证据几乎同时抵达：1905 年爱因斯坦用布朗运动（第 25 章提到的花粉跳舞）算出水分子的“撞击账”，佩兰随后用显微镜下的实验证实了它——原子和分子真实存在，统计解释完全正确。今天，$S=k\ln W$ 刻在玻尔兹曼的墓碑上。这个故事值得记住的不是胜负，而是方法：一个理论只要给出可检验的预言（比如布朗运动的精确统计规律），再激烈的哲学反对也终将被实验裁决。
\end{zhengju}

\section{这套账本的边界}

$\Delta G$ 这杆秤极好用，但它的使用说明书上有几条小字：

\begin{itemize}[itemsep=2pt]
\item \textbf{只判恒温恒压下的方向}。深空、地心、高压反应釜里，条件变了，账要重算（还有别的“自由能”管那些场合）。
\item \textbf{只判方向，不判快慢}。如前所述，$\Delta G<0$ 的过程可以慢到等于不发生；反过来，想让它快起来，得研究活化能和催化剂（第 29、30 章）。
\item \textbf{判的是“能不能”，不是“动不动”}。$\Delta G<0$ 的体系可能停在路上；$\Delta G=0$ 的体系也没“静止”，只是正逆两边扯平（第 31 章会看到，平衡其实是两边都在跑）。
\item \textbf{熵增的铁律只管孤立体系}。你的房间、你的身体、地球都不是孤立体系——地球一直在接收太阳的“低熵”阳光、向太空排出“高熵”的热辐射。局部熵减（结冰、生长、盖起一座城市）只要配上别处更大的熵增，就合法。拿“熵增”断言“宇宙里不该出现有序结构”，是用错了定律的适用范围。
\end{itemize}

\section{再看那只冰袋}

回到体育课的医药箱。现在你手里有完整的账：硝酸铵溶解，断开了晶格里的离子吸引、又要重建水合外壳，总账吸热，$\Delta H>0$——这就是袋子变冷的原因，冷正是“亏损”从你皮肤上收来的。但晶体瓦解成自由的水合离子，$\Delta S$ 大增；室温的 $T$ 乘上这笔熵增，压过了焓的亏损，$\Delta G<0$，于是溶解自发进行，一路进行到离子浓度高得让账扯平为止。

实验里那三个“却”，现在逐条对上：它\textbf{自发进行}，因为 $\Delta G<0$，不需要任何人推；它\textbf{自发还吸热}，因为熵账赢过了焓账，温度越高赢得越稳；它\textbf{外壁始终不结冰}，因为一袋硝酸铵能从环境抢走的热量有限，袋子温度降上十几度就到底了，离 $0\,^{\circ}\mathrm{C}$ 还远——而就算真到了零度以下，溶满离子的水也比纯水更难结冰（稀溶液的这个脾气，第 25 章的挑战层算过账）。

它也把三句容易混的话一次讲清：它自发，却\textbf{不放热}（焓不是唯一裁判）；它自发，却\textbf{有尽头}（$\Delta G$ 随浓度变化，趋向零）；它自发，而且恰好很快——但快是它的运气，不是“自发”的赠品，金刚石就没这份运气。

\section{方向之后，轮到速度}

这一章我们给化学反应办齐了“方向证”：$\Delta H$ 和 $\Delta S$ 两本账经温度折汇率，合并成 $\Delta G$，正负定方向，温度能翻盘。这套账本威力极大——冶金工程师靠它算出“多少度才能把矿石里的铁夺回来”，化工工程师靠它判断“这个反应有没有做的必要”，医生也离不开它：疫苗和血浆冻干保存，靠的就是降温降压把“冰变水”的方向整个扳掉，让水直接以气体形式离开，样品既不腐败也不被烫坏。

但办齐了方向证，只走完一半。金刚石守着“该变石墨”的判决书几亿年不动身；炸药和铁锈同为氧化，快慢差着十几个数量级；夏天的食物一夜之间馊掉，冰箱里的能多撑一周。方向归热力学管，快慢归另一套规律管——粒子要撞得多狠、翻过多高的坎，才能真的换身。下一章（第 29 章），我们就去量那道坎的高度。

\begin{tiaozhan}
来亲手算一笔玻尔兹曼的账。他的公式是 $S=k\ln W$，其中 $W$ 是微观排法数，$k$ 是玻尔兹曼常数，\ud{1.38\times10^{-23}}{J/K}。取 \ud{1}{mol} 气体（$6.02\times10^{23}$ 个分子），让它从盒子的一半扩散到整个盒子：每个分子的“可选位置”翻倍，$W$ 就变成原来的 $2^{6.02\times10^{23}}$ 倍。于是
\[\Delta S = k\ln\left(2^{6.02\times10^{23}}\right) = k\times 6.02\times10^{23}\times\ln 2 \approx 8.31\times 0.693 \approx \ud{5.8}{J/K}\]
注意那个 $8.31$：它就是第 26 章账本里常见的摩尔气体常数 $R$。宏观世界的 $R$ 和微观世界的 $k$，原来隔着一座“\ud{6.02\times10^{23}}{}”的桥——这正是统计力学最漂亮的风景：热力学的每个宏观常数，背后都站着一次微观计数。跳过这段不影响主线，但你若算通了，就亲手摸到了玻尔兹曼墓碑上那行公式的温度。
\end{tiaozhan}

\section*{练一练}

\begin{sikaoti}
\item 画一幅“冰袋粒子图”：左边画溶解前（整齐的硝酸铵晶格加独立的水分子），右边画溶解后（分散的水合离子）。用箭头标出热量的流向，并分别注明体系和环境的熵是增还是减。图旁注明：这是人为的表示法，不代表真实外观。
\item 不用公式，只用“排法多少”解释：为什么一滴红墨水滴进清水，最终会均匀散开，而你永远不会看到散开的墨水自己聚回一滴？如果拍成视频倒放，倒放版违反了哪条力学定律吗？
\item 水结冰放热（$\Delta H<0$）、熵减（$\Delta S<0$）。用 $\Delta G=\Delta H-T\Delta S$ 解释：为什么冰箱冷冻室 $-18\,^{\circ}\mathrm{C}$ 里水一定结冰，而 $10\,^{\circ}\mathrm{C}$ 的冷藏室里绝不结冰？冰与水恰好“扯平”的温度满足什么条件？
\item 把下列过程按“体系熵增/熵减”分类，并各用一句话说明理由：香水挥发；水蒸气在镜子上凝成雾；小苏打受热放出 \chem{CO_2}；植物把 \chem{CO_2} 和水组装成葡萄糖。
\item “纸是可燃的，燃烧 $\Delta G<0$，所以这本书随时可能自己烧起来。”这句话错在哪一步？请把“方向”和“速度”分开反驳，并指出真正拦住纸张的那道坎叫什么（提示：下一章的主角）。
\item 查资料找出石灰石煅烧的温度，结合正文中“升温使吸热且熵增的反应转自发”的道理，向家人解释：为什么石灰窑要烧那么高的温度，而不是“多等一会儿”让反应在室温慢慢进行。
\end{sikaoti}
